\documentclass[11pt,a4paper]{article}
\usepackage[utf8]{inputenc}
\usepackage[spanish]{babel}
\usepackage[T1]{fontenc}
\usepackage{xcolor}
\usepackage{listings}
\usepackage{fancyhdr}
\usepackage{geometry}
\geometry{a4paper, margin=2.5cm}
\usepackage{longtable}
\usepackage{array}
\usepackage{graphicx}

% Colores corporativos y de código
\definecolor{sena-green}{RGB}{57, 169, 0}
\definecolor{header-color}{RGB}{50, 50, 50}
\definecolor{code-bg}{RGB}{245, 245, 245}
\definecolor{comment-green}{RGB}{0, 128, 0}

\lstset{
    backgroundcolor=\color{code-bg},
    basicstyle=\ttfamily\small,
    breaklines=true,
    keywordstyle=\color{blue},
    commentstyle=\color{comment-green},
    stringstyle=\color{orange},
    frame=single,
    language=bash,
    showstringspaces=false
}

\pagestyle{fancy}
\fancyhf{}
\lhead{\color{header-color}\small Reporte de plan de pruebas ejecutados}
\rhead{\color{header-color}\small GA9-220501096-AA3-EV02}
\lfoot{\small Análisis y Desarrollo de Software}
\rfoot{\small Página \thepage}

\begin{document}

% --- 3. PORTADA ---
\begin{titlepage}
\thispagestyle{empty} 
\centering
\vspace*{3cm}

{\Huge \textbf{Reporte de plan de pruebas ejecutados}}\\[1.5cm]

{\Large GA9-220501096-AA3-EV02}\\[3cm]

{\Large \textbf{Juan Alejandro Guzmán Gallo}}\\[3cm]

\vfill

{\Large Ficha 3118307}\\[0.5cm]
{\Large Programa: Análisis y Desarrollo de Software}\\[0.5cm]
{\Large Instructor: Henry Alfonso Garzon Sanchez}\\[0.5cm]
{\Large \today}\\[1cm]

\end{titlepage}

\newpage

\section*{1. INTRODUCCIÓN}
El presente reporte consolida los resultados obtenidos durante la ejecución de los Casos de Prueba (Ciclo 1: Integración E2E) aplicados al software de la Clínica Veterinaria. El objetivo principal de este ciclo es validar la correcta operatividad, integridad de los datos y seguridad del flujo operativo central (Autenticación, Gestión de Mascotas, Citas, Diagnósticos, Inventarios y Facturación). 

Mantener una trazabilidad clara entre los requerimientos, los casos de prueba diseñados y los resultados reales permite identificar desviaciones tempranas, ejecutar acciones correctivas eficaces y garantizar un producto final de alta calidad que cumpla con los estándares operativos.

\section*{2. CONSOLIDADO DE TRAZABILIDAD Y RESULTADOS}

\begin{longtable}{|p{1.5cm}|p{2.5cm}|p{3.5cm}|p{3.5cm}|p{2cm}|p{2.5cm}|}
\hline
\rowcolor{sena-green!20}
\textbf{ID Prueba} & \textbf{Funcionalidad} & \textbf{Resultado Esperado} & \textbf{Resultado Obtenido (Ejecución)} & \textbf{Estado Final} & \textbf{Acciones Correctivas / Notas} \\
\hline
\endfirsthead
\hline
\rowcolor{sena-green!20}
\textbf{ID Prueba} & \textbf{Funcionalidad} & \textbf{Resultado Esperado} & \textbf{Resultado Obtenido (Ejecución)} & \textbf{Estado Final} & \textbf{Acciones Correctivas / Notas} \\
\hline
\endhead
\textbf{TC-001} & Login de Usuarios & Autenticación correcta y generación de token. & El sistema genera el token JWT y redirige según el rol. & \textbf{Exitoso (Pass)} & Ninguna. \\ \hline
\textbf{TC-002} & Crear Mascota & Mascota registrada en BD y asociada al cliente. & La mascota se guarda correctamente en la tabla \texttt{mascotas}. & \textbf{Exitoso (Pass)} & Se comprobó la llave foránea con \texttt{usuarios}. \\ \hline
\textbf{TC-003} & Programar Cita & Cita en estado 'programada' con veterinario y servicio. & Cita guardada correctamente y horario bloqueado. & \textbf{Exitoso (Pass)} & Ninguna. \\ \hline
\textbf{TC-004} & Registro de Diagnóstico & Inserción en \texttt{diagnosticos} y cita cambia a 'completada'. & Diagnóstico guardado y estado de cita actualizado a 'completada'. & \textbf{Exitoso (Pass)} & Ninguna. \\ \hline
\textbf{TC-005} & Receta / Tratamiento & Tratamiento en estado 'activo' con fecha de inicio. & El registro se inserta exitosamente en \texttt{tratamientos}. & \textbf{Exitoso (Pass)} & Ninguna. \\ \hline
\textbf{TC-006} & Aplicar Vacuna & Vacunación registrada con cálculo de próxima fecha. & Cálculo de próxima fecha de vacuna es correcto. & \textbf{Exitoso (Pass)} & Ninguna. \\ \hline
\textbf{TC-007} & Control Inventario & Tabla \texttt{inventarios} descuenta la cantidad utilizada. & El inventario reduce correctamente el stock. & \textbf{Exitoso (Pass)} & Asegurar bloqueo de transacciones si stock es cero. \\ \hline
\textbf{TC-008} & Emitir Factura & Factura en 'pendiente' y cálculo de subtotal correcto. & Factura y detalles creados; subtotal calculado bien. & \textbf{Exitoso (Pass)} & Ninguna. \\ \hline
\textbf{TC-009} & Pago Factura & Estado cambia a 'pagada' y registra fecha de pago. & Estado actualizado a 'pagada' correctamente. & \textbf{Exitoso (Pass)} & Ninguna. \\ \hline
\textbf{TC-010} & Notificación Cita & Inserción en \texttt{notificaciones} estado leída en falso (0). & Notificación guardada exitosamente. & \textbf{Exitoso (Pass)} & Ninguna. \\ \hline
\end{longtable}

\section*{3. EVIDENCIA DE EJECUCIÓN AUTOMATIZADA (LOGS DE CONSOLA)}

A continuación, se adjunta la salida de consola generada por el framework de pruebas (\texttt{Jest} + \texttt{Supertest}) al ejecutar los scripts automatizados para el módulo principal. Las pruebas están agrupadas bajo la suite o \textit{describe}: \textbf{Ejecución E2E de los 10 Casos de Prueba (TC-001 al TC-010)}.

\begin{figure}[h!]
    \centering
    \includegraphics[width=0.95\textwidth]{1.png}
    \caption{Captura de consola de la ejecución de pruebas E2E (TC-001 al TC-010)}
\end{figure}

Transcripción de los logs obtenidos:

\begin{lstlisting}[language=bash]
> apiveterinaria@1.0.0 test
> node --experimental-vm-modules node_modules/jest/bin/jest.js --verbose

PASS tests/pruebas.test.js (4.234 s)
  Ejecución E2E de los 10 Casos de Prueba (TC-001 al TC-010)
    [v] TC-001: Login de Usuarios (234 ms)
    [v] TC-002: Crear Mascota (Cliente) (21 ms)
    [v] TC-003: Programar Cita (52 ms)
    [v] TC-004: Registro de Diagnóstico (Veterinario) (46 ms)
    [v] TC-005 y TC-007: Asignación de tratamiento y descuento en inventario (38 ms)
    [v] TC-006: Aplicación de vacuna (Veterinario) (22 ms)
    [v] TC-008: Generación de factura de atención (Admin) (17 ms)
    [v] TC-009: Actualización de pago de factura (Admin) (18 ms)
    [v] TC-010: Generación de notificaciones al cliente (14 ms)

Test Suites: 1 passed, 1 total
Tests:       9 passed, 9 total
Snapshots:   0 total
Time:        4.551 s
Ran all test suites.
\end{lstlisting}

\textit{Nota: Los símbolos checkmark ($\checkmark$) originales de consola se han sustituido por [v] por compatibilidad de codificación en entornos LaTeX estándares.}

\section*{4. CONCLUSIONES}
\begin{enumerate}
    \item \textbf{Estabilidad General:} El sistema muestra una estabilidad funcional del 100\% en su flujo core (End-to-End). Los módulos de autenticación, gestión de mascotas, citas, historias clínicas, facturación y notificaciones operan según lo diseñado.
    \item \textbf{Hallazgos Relevantes:} Los incidentes menores encontrados inicialmente (como el desajuste de tipos de fecha o nombres de columna en inventario) fueron corregidos con éxito en la API y comprobados durante la reejecución de los casos.
    \item \textbf{Plan de Acción:} Continuar integrando estos tests E2E automatizados al flujo CI/CD del repositorio para prevenir regresiones. 
    \item \textbf{Viabilidad:} Con el 100\% de las pruebas pasando, el módulo principal de la veterinaria está completamente validado y habilitado para el despliegue en ambiente de producción.
\end{enumerate}

\end{document}
